\chapter{本地量化推理引擎}

本章把全书的线索收束到贯穿项目：Linux CPU Quantized Inference。我们要在本地 Linux 上加载一个小型量化模型，执行推理，验证输出，并记录性能。这个项目规模不大，但它包含推理引擎的核心骨架，后续可以继续扩展到卷积、Attention 和更复杂模型格式。

\section{模型文件}

教学模型文件采用文本格式，便于手工检查。文件记录输入维度、隐藏维度、输出维度、每层 scale、int8 权重和 float bias。真实工业模型常用二进制格式，是为了加载速度、体积和元数据完整性；教学阶段先用文本格式，是为了让读者看清每个数字进入推理路径的位置。

\section{量化线性层}

第一阶段使用 int8 权重和 float 激活。计算时把权重按 scale 反量化，再与输入相乘累加。这个实现不是最高效的 int8 GEMM，但语义清楚，适合作为正确性基线。后续可以把反量化和累加融合，或使用 int32 累加和输出阶段 scale。

\begin{lstlisting}[language=C++,caption={量化线性层的第一阶段语义}]
float weighted_sum_i8(const std::int8_t* weights,
                      const float* input,
                      float scale,
                      std::int32_t input_size) {
    float acc = 0.0F;
    for (std::int32_t i = 0; i < input_size; ++i) {
        acc += static_cast<float>(weights[i]) * scale * input[i];
    }
    return acc;
}
\end{lstlisting}

\section{推理执行}

推理执行器先运行隐藏层，写入隐藏激活并执行 ReLU；再运行输出层，得到 logits；最后取最大 logit 的下标作为分类结果。这个流程简单，但已经有模型参数、激活缓冲区、算子调度和输出解释。测试用固定输入和模型文件验证结果，benchmark 重复执行多次记录平均延迟。

\section{正确性和性能报告}

正确性报告至少包含模型文件、输入、期望输出、实际输出和误差。性能报告至少包含 CPU 信息、编译器、编译选项、线程数、重复次数、平均延迟和吞吐。若开启多线程，还要记录任务切分方式。若后续引入 SIMD 或量化累加优化，要保留标量基线对比。

\begin{keyidea}
推理引擎不是一个巨大函数，而是一组边界清晰的模块：模型加载、张量表示、量化权重、算子内核、执行器、测试和 benchmark。模块边界清楚，后续优化才不会破坏正确性。
\end{keyidea}

\section{后续扩展}

LCQI 的自然扩展路径有四条。第一，支持更多算子，例如卷积、RMSNorm、Softmax 和 Attention。第二，改进量化格式，例如 per-channel scale、group-wise quantization 和 int32 累加。第三，加入线程池和更稳定的任务切分。第四，支持更接近真实模型的二进制模型格式。每次扩展都必须保留测试和 benchmark。

\begin{labbox}
本章实验：构建 LCQI，运行测试和 benchmark。确认模型文件可以在本地 Linux 上加载，推理输出分类下标正确，并记录单次推理平均耗时。
\end{labbox}
